\vspace{-1em}
\begin{abstract}
    While talking head generation has advanced rapidly, generating natural listener behavior in dyadic conversations, which know when to react, how to react, and with what type of response, remains underexplored. Existing dyadic datasets lack fine-grained listener reaction annotations, and prevailing evaluation metrics inherited from talking-head and video generation measure visual realism rather than whether a listener reacted appropriately. We address these gaps along three aspects. First, we curate a listening-head-specific dataset built from RealTalk and Seamless Interaction, comprising approximately 147 hours of paired speaker–listener videos with 64,557 event-level reaction annotations across six categories: \textit{nodding}, \textit{head shaking}, \textit{smiling}, \textit{laughing}, \textit{frowning}, and \textit{surprised}. Second, we introduce an audio-driven baseline built on a flow-matching transformer, namely GLARE, with prosody conditioning derived from Qwen2-Audio and a temporal reaction loss that explicitly supervises frame-wise reactions. Third, we propose a reaction-oriented evaluation protocol that jointly measures reaction occurrence (R-F1), temporal alignment (R-tIoU), asymmetric temporal deviation (R-ATD), and reaction-region visual quality (R-FID), giving a more behaviorally grounded assessment than visual-quality-only metrics. Experiment results show consistent gains over prior listening-head methods in both visual fidelity and reaction-level metrics, suggesting that reaction-aware data, modeling, and evaluation are critical for natural listening behavior. The code, annotation pipeline, and processed dataset will be released at \href{https://github.com/lzk901372/glare}{https://github.com/lzk901372/glare}.
\end{abstract}